%----------------------------------------------------------
% [표 예시 1] tabularray(tblr): 열 너비 자동 조정, 세로 병합
%  - colspec: Q = 내용 폭, X = 남은 폭을 채움 / [c,m] = 가운데·세로 중앙
%  - cell{2}{3} = {r=2}{} : 2행 3열 셀을 아래로 2칸 병합(병합되는 칸은 비워 둠)
%----------------------------------------------------------
\begin{table}[!htbp]
\centering
\small
\caption{Summary of datasets (example of a \texttt{tblr} table).}
\label{tab:dataset}
\begin{tblr}{
  width = \linewidth,
  colspec = { Q[c,m] X[l,m] Q[c,m] Q[c,m] },
  row{1} = {font=\bfseries, halign=c},
  cell{2}{3} = {r=2}{},
}
\toprule
Dataset   & Description                                   & Type          & Resolution \\
\midrule
Dataset A & Description of dataset A                      & {Satellite \\ imagery} & 10 m \\
\addlinespace
Dataset B & Description of dataset B                      &               & 30 m \\
\addlinespace
Dataset C & Description of dataset C (e.g., in-situ data) & Point data    & -- \\
\bottomrule
\end{tblr}
\end{table}
